% sec-00-front.tex - Abstract, reader's guide, and the two indexes.  FINAL stage.
% Senior-author pass.  The abstract is carried forward unchanged at 1347
% characters; the master prompt fixes that budget and no stage may exceed it.
% Column widths in Tables 1, 2 and 3 are re-cut against the compiled widths
% rather than the estimated ones: the Platform column of Table 2 loses 0.10 cm
% and the Contents column of Table 3 gains it, because Table 2's longest
% platform string is nine characters and Table 3's longest contents string is
% fifty-two.  Every caption is checked to two lines within a four-character
% spread.  No drafting instruction survives in this file.
\section*{Abstract}
\phantomsection\addcontentsline{toc}{section}{Abstract}

The prior pancreatic cancer trial system built its documents on a month
scale, through teams, with peer review arriving after completion. The new
system produces the same documents on a 1 to 4 day project time scale, through
one author directing a master prompt, with AI peer review arriving during
development. This paper discloses that system across four deliverables. The IND
assembles twelve modules under 21 CFR 312 in four days. Two protocols follow: a
Phase 1 first-in-human study at n = 18 using a 3+3 escalation behind a Phase 0
gate of at least 1000 simulations across two frameworks, and a Phase 2
randomized study at n = 220 across eight centers targeting a hazard ratio of
0.60. Five bill versions and four companion documents were deposited in eleven
days. Fourteen funding artifacts followed, including a \$1,606,000
small-business ask against a \$3,500,000 five-year direct program, and a virtual
trial projected at \$36,330 against an industry benchmark above \$120,000.
Human peer review cannot absorb this output. A best-case round is seven to eight
weeks and a typical one several months, against three model manufacturers
reporting the same day; and review arriving after release cannot correct the
released object. The two systems are incompatible, and the comprehensive
provenance record the new one leaves is why it is trusted.

\section*{How to Read This Document}
\phantomsection\addcontentsline{toc}{section}{How to Read This Document}

This paper is a disclosure of a working method, and its evidence is a set of
artifacts that already exist and are already deposited. It is written so that
four different readers can each open it at a different page and find the thing
they came for within two pages. Table 1 names the four and the section each
should open first. No reader is asked to read the document front to back, and
every entry point named in Table 1 is a table or a figure rather than a
paragraph.

\begin{apptable}
\begin{tabularx}{\textwidth}{>{\raggedright\arraybackslash}p{2.85cm} >{\raggedright\arraybackslash}p{2.35cm} >{\raggedright\arraybackslash}p{2.35cm} >{\raggedright\arraybackslash}X}
\headrow \thead{Reader} & \thead{Open at} & \thead{Then} & \thead{The question this reader is actually asking}\\
\midrule
Funding officer & \S\ref{sec:funding}, Table 18 & \S\ref{sec:methods}, Figure 3 & What was produced, at what cost, and by what repeatable method? \\
Regulatory reviewer & \S\ref{sec:ind}, Table 9 & \S\ref{sec:protocol}, Figure 10 & Does the dossier map onto the codified content items, and where? \\
Legislative staffer & \S\ref{sec:legislation}, Table 14 & \S\ref{sec:legislation}, Figure 14 & What duty does the bill create, and who is positioned to discharge it? \\
Peer reviewer & \S\ref{sec:peerreview}, Table 21 & \S\ref{sec:peerreview}, Figure 24 & How was this reviewed, and what happened when the reviewers disagreed? \\
\end{tabularx}
\end{apptable}
\vspace{-0.6cm}
\tabcap{Table 1. Four readers, the section each should open first, and the one\\question
each is actually asking of this paper's method and record.}

\vspace{0.25cm}

Two indexes follow, and they are the paper's other entry points. Table 2 lists
the twenty-five figures with the platform each is drawn on and the one question
each answers, and Table 3 does the same for the twenty-five tables. A reader who
knows the question but not the page should start at whichever of the two carries
it. The figure count per platform is set by purpose and not by quota: six
mermaid figures because six claims are chronological, six d2 figures because six
are grids or containers, five graphviz because five are structural, four
plantuml because four need a guard or a fork to be stated honestly, and four
clustered infrastructure figures because four are about a system's parts.

\begin{apptable}
\begin{tabularx}{\textwidth}{>{\raggedright\arraybackslash}p{0.45cm} >{\raggedright\arraybackslash}p{0.4cm} >{\raggedright\arraybackslash}p{1.65cm} >{\raggedright\arraybackslash}p{2.85cm} >{\raggedright\arraybackslash}X}
\headrow \thead{Fig} & \thead{\S} & \thead{Platform} & \thead{Source directory} & \thead{The question it answers}\\
\midrule
1 & 1 & Mermaid & \appfile{mermaid/} & \footnotesize Which capability do eleven Federal actions leave unfilled? \\
2 & 1 & D2 & \appfile{d2/} & \footnotesize On which ten axes do the prior and new systems differ? \\
3 & 2 & PlantUML & \appfile{plantuml/} & \footnotesize Which part of the build is concurrent, and which strictly serial? \\
4 & 2 & Mermaid & \appfile{mermaid/} & \footnotesize What happens between one prompt and one pushed commit? \\
5 & 2 & Graphviz & \appfile{graphviz/} & \footnotesize What does a figure specification store that a raster cannot? \\
6 & 3 & Diagrams & \appfile{diagrams-python/} & \footnotesize Which source clusters fed which IND module? \\
7 & 3 & Mermaid & \appfile{mermaid/} & \footnotesize How long did twelve IND modules take, against the prior system? \\
8 & 3 & D2 & \appfile{d2/} & \footnotesize Which file satisfies each codified IND content item? \\
9 & 3 & Graphviz & \appfile{graphviz/} & \footnotesize What combination of failures would produce a clinical hold? \\
10 & 4 & PlantUML & \appfile{plantuml/} & \footnotesize What guard fires each transition in a participant's state machine? \\
11 & 4 & Mermaid & \appfile{mermaid/} & \footnotesize What must hold on each rung from Phase 0 to Phase 2? \\
12 & 4 & D2 & \appfile{d2/} & \footnotesize What did Phase 2 carry, replace, and add? \\
13 & 4 & Diagrams & \appfile{diagrams-python/} & \footnotesize What is inside the site boundary, and what is outside it? \\
14 & 5 & PlantUML & \appfile{plantuml/} & \footnotesize Which actor owns each statutory duty? \\
15 & 5 & Graphviz & \appfile{graphviz/} & \footnotesize What did each bill version add that its predecessor lacked? \\
16 & 5 & D2 & \appfile{d2/} & \footnotesize How does one requirement reach the operating room? \\
17 & 6 & Mermaid & \appfile{mermaid/} & \footnotesize When did each funding artifact appear, and how fast? \\
18 & 6 & D2 & \appfile{d2/} & \footnotesize What does the same direct work cost under two overhead regimes? \\
19 & 6 & Graphviz & \appfile{graphviz/} & \footnotesize Where does each dollar land, and what does the award not reach? \\
20 & 6 & Diagrams & \appfile{diagrams-python/} & \footnotesize What machinery turned one evidence store into fourteen artifacts? \\
21 & 7 & Mermaid & \appfile{mermaid/} & \footnotesize How long does one review round take under each regime? \\
22 & 7 & D2 & \appfile{d2/} & \footnotesize What is the ratio between the two regimes on six economic axes? \\
23 & 7 & PlantUML & \appfile{plantuml/} & \footnotesize Who reviews concurrently, and where does a human gate sit? \\
24 & 7 & Graphviz & \appfile{graphviz/} & \footnotesize What happens when three model reviewers disagree? \\
25 & 8 & Diagrams & \appfile{diagrams-python/} & \footnotesize What exposures remain, and what mitigates each? \\
\end{tabularx}
\end{apptable}
\vspace{-0.6cm}
\tabcap{Table 2. The twenty-five figures, their section, their platform, the\\
directory carrying each one, and the question that each of them answers.}

\begin{apptable}
\begin{tabularx}{\textwidth}{>{\raggedright\arraybackslash}p{0.6cm} >{\raggedright\arraybackslash}p{0.4cm} >{\raggedright\arraybackslash}p{6.0cm} >{\raggedright\arraybackslash}X}
\headrow \thead{Tab} & \thead{\S} & \thead{Contents} & \thead{The question it answers}\\
\midrule
1 & 0 & Four readers, four entry points & Which section should this reader open first? \\
2 & 0 & Figure index, twenty-five rows & Where is each figure and what does it answer? \\
3 & 0 & Table index, twenty-five rows & Where is each table and what does it answer? \\
4 & 1 & Federal AI and cancer actions, Jan 2025 to Jul 2026 & What has the government already committed to? \\
5 & 1 & Prior system against new system, ten axes & What exactly changed, measured? \\
6 & 2 & Master prompt and sub-prompt counts across works & How many times has this method been run? \\
7 & 2 & Machine-readable diagram storage and reuse & What does the specification format buy? \\
8 & 3 & IND package inventory by module & What is in the dossier? \\
9 & 3 & IND section, source, and generation window & Where did each module come from, and how fast? \\
10 & 3 & IND complexity and volume record & What scale did one author reach? \\
11 & 4 & Phase 1 protocol synopsis & What is the first-in-human study? \\
12 & 4 & Phase 2 protocol synopsis and deltas & What changed when it randomized? \\
13 & 4 & Protocol generation metrics, both protocols & How long did two protocols take? \\
14 & 5 & Statutory sections and the duty each creates & What does the bill require of whom? \\
15 & 5 & Bill version lineage, v1.0 to v5.0 & What did each version add? \\
16 & 5 & Legislative artifact metrics & How much statutory text, how fast? \\
17 & 6 & Ten funding applications and mechanism class & Who was asked, and under what authority? \\
18 & 6 & Capitalization numbers & What is the ask, and what does it buy? \\
19 & 6 & RFA-RM-27-001 v2 structure & What does a full NIH application contain? \\
20 & 6 & Cost and time benchmarks, prior against new & What does the same work cost each way? \\
21 & 7 & Prior against new peer review, side by side & Why can the prior system not absorb this? \\
22 & 7 & Model roster and review role & Which model does what, and for whom? \\
23 & 7 & AI peer review quality and efficiency metrics & How was review quality measured? \\
24 & 7 & AI peer review instances across author works & How often has this been done before? \\
25 & 8 & Limitations register with mitigations & What is still exposed, and what reduces it? \\
\end{tabularx}
\end{apptable}
\vspace{-0.6cm}
\tabcap{Table 3. The twenty-five tables, their section, their contents, and the one\\question
each answers, so a reader who knows the question can find the page.}

\section*{How This Paper Was Built}
\phantomsection\addcontentsline{toc}{section}{How This Paper Was Built}

One master prompt was supplied by the author on August 14, 2026. Claude Code
Opus 5 decomposed it into eight sub-prompts and executed them in order without
further author intervention, committing and pushing every generated file the
moment it was finished. The eight stages are recorded in
\appfile{new-trial-system/sub-prompts}, and the master prompt itself is filed
verbatim in \appfile{new-trial-system/prompts/prompt-new-trial.md}.

Stages 1 through 5 produced twenty-five machine-readable figure specifications,
one file per figure, across five diagram platforms. Each specification carries a
perspective statement naming what no other figure in the paper shows, the
figure's two-line caption exactly as printed here, valid source in the
platform's own syntax, a TikZ construction table of absolute coordinates with a
stated node pitch, an edge-routing paragraph naming every edge that could cross
a node and the clearance that prevents it, a value table attributing every
number in the figure to a repository file, and a list of the exact sources read.
No figure in this paper was drawn from memory of another figure, and no figure
is a raster.

Stages 6, 7 and 8 produced the paper three times. The draft fixed the section
set, the figure numbering and the table numbering, and carried a bracketed
instruction in every section naming the repository file the next stage had to
read. The full version drew all twenty-five figures, populated all twenty-five
tables, discharged every instruction, and wrote the prose from the four
author-final \appfile{publication/} directories and the four input archives with
direct quotation. This final version is the senior author's proof-reading pass:
\S\ref{sec:introduction} through \S\ref{sec:conclusions} carry the same argument
as the full version, set correctly on the page.

Each stage produced a self-contained Overleaf bundle, and each bundle is
deposited beside its source. A reader who wants to check the production-rate
claim of \S\ref{sec:methods} against the record can read the branch history
rather than this paragraph.
